\documentclass{article}
\usepackage{amsmath,amssymb,hyperref}
\title{Global joining comparisons for finite-stage evaluators}
\date{15 November 2026}
\begin{document}\maketitle
\paragraph{Status.} Constructive conditional theorem with a written proof. The finite-stage evaluators $F_n:T_n\to D$ are available under the independently specified higher-uniformity and assembly doctrine. This note does not assume their compatibility with ordered joining. It constructs the comparisons from additional semantic joining data and supplies the transported unit-frame normalization for its final coherence.

\paragraph{Target joining input.} Use pseudonatural adjoint equivalences
$j_{X,Y,f}:B_XB_Y(f)\simeq B_{Y\otimes X}(\bar f)$ and $v:B_I\simeq1$, including the actual action transports. Require the full two-boundary native three-cell comparison, natural in each admitted witness $(\sigma,\Gamma)$ and compatible with pasting such rectangles. Require joining's specified sequential and parallel assembly three-cell comparisons, natural on the same full witnesses and satisfying their unit/composition equations. Triple-joining associativity and joining-unit comparisons are invertible target three-cells; their four-closure pentagon and unit diagrams are equations. These are semantic target data, not consequences of an ordinary feedback function.

\paragraph{Fix the direction and the reduction.} For monotone $\theta:n\to m$ we seek
\[
 E_\theta:F_n\Longrightarrow F_mR_\theta.
\]
For a two-coordinate merge this uses forward $j$; for an empty-fiber insertion it uses $v^{-1}$. This is the adjoint direction to a comparison $F_mR_\theta\Rightarrow F_n$. At a machine, first use the actual raw/normalized and action comparisons to display $F_n$ as the nested closure word $B_{X_1}\cdots B_{X_n}(F_0f)$, with its specified body parentheses. In each nonempty fiber reduce the consecutive block by merging its first two outer-to-inner closures, then merging that result with the next closure, retaining all reassociations of the reversed total state. Process the fibers in descending target-index order. Insert the empty-fiber $B_I$ using $v^{-1}$ in those fixed positions. Finally use the specified action comparisons to identify the resulting block word with $F_mR_\theta$. This defines one definite paste; path independence is not used as its definition.

\paragraph{Native tritransformation data.} The interface component of $E_\theta$ is $1_A$ in $D$, with the actual left/right identity whiskers. Its component at a machine is the unitor-conjugated two-cell just constructed. At a native arrow, paste the two-boundary three-cell for each local $j$ reduction and the naturality cell for each $v^{-1}$ insertion. This includes the images of all inner state maps and actual process witnesses, not just a bare map of the outer state. Naturality at a full source three-cell is the corresponding stipulated equality in target uniformity. Pasting and identity naturality follow from the target rectangle-pasting laws. The composition and identity compatibility modifications use the supplied joining/assembly comparisons, with the same unit frames. These statements are precisely the data and local equations to check for a tritransformation; they do not infer it from the machine component alone.

\paragraph{Reduction-substitution lemma.} Compare any composite of the fixed fiber reductions with the fixed reduction of its total ordered partition. Reduce its finite parenthesization tree to the chosen comb using the target triple-joining comparison; remove inserted unit nodes using the target joining-unit comparison. For a reduction occurring inside a larger closure expression, whisker that comparison by the surrounding $B$ pseudofunctors and their explicit composition constraints. Naturality of the joining comparisons identifies this contextual reduction with the globally framed reduction. Two overlapping rotations are governed by the target four-closure pentagon, rotation/unit overlaps by the joining-unit equations, and disjoint reductions by pseudofunctor interchange and naturality. Induction on the surrounding tree establishes substitution compatibility. This is a finite formal-word argument with the actual typed target comparisons; it does not invoke faithfulness of a state projection of $D$.

\paragraph{Pair and triple indexing coherence.} For $\theta,\phi$, compare the two tritransformation pastes from $F_n$ to $F_lR_\phi R_\theta$ (after the evaluator image of the specified source composition comparison $c_{\phi,\theta}$). The reduction-substitution lemma supplies the machine three-cells of a trimodification; its interface component uses the transported unit-frame normalization defined below, with the chosen whiskering frames. Full-witness naturality and assembly compatibility follow by the same substitution applied to the given target laws.

For three indexing arrows, compatibility of these trimodifications with the source joining trimodification $M$ requires a perturbation. Choose its interface components by transporting the admitted source unit-frame comparisons through the normalization below, rather than replacing them by unspecified canonical target cells. Its equation at a machine is the reduction-substitution equation for three stages of fiber flattening, hence the target joining pentagon and unit equations with the source comparison frames inserted. Admissibility must be checked before declaring this a perturbation. No three-cell is assumed unique merely because its state shadow would be unique.

\paragraph{The full source unit-frame fragment.} At an external interface retain every formal unit-state word arising in the actual comparison pastes, including uneven tuples $F_\theta(I^n)$, the unit and binary monoidal comparisons of $F_\theta$ on such words, and their unit braids. The binary comparison can occur through whiskering; it is not excluded. For each word $W$ in $S^m$ choose a coordinatewise normalization $\nu_W:W\simeq I^m$ with its admitted pure process square. For example $(I\otimes I,I)$ uses $(\lambda_I,1_I)$; it is not literally a tensor word in the sole object $I_{S^2}$.

For each structural map $q:W\to V$ in this fragment choose the structural invertible two-cell
\[
 u_q:\nu_Vq\Rightarrow\nu_W.
\]
Both sides have the same formal empty-word permutation. Structural symmetric monoidal coherence supplies $u_q$ and its identity, composition and structural-two-cell equations. Its native admission uses the actual pure action/unit body comparisons. Hence this is a coherently contractible \emph{chosen structural fragment}; no assertion is made about arbitrary state endomorphisms of the unit. The relevant source is Gurski--Osorno, \href{https://arxiv.org/pdf/1210.1174v1}{Theorem 1.8 and Corollary 1.9}. Its use is only on the residual formal structural diagram after admission.

\paragraph{A local transport lemma.} The normalization step is a bicategorical fact, not a faithfulness argument. Let $\mathcal C$ be a chosen sub-bicategory of a source hom bicategory, with object $i$ and adjoint equivalences $\nu_W:W\simeq i$. Suppose the invertible cells $u_q:\nu_Vq\Rightarrow\nu_W$ satisfy identity, composition and two-cell naturality, so they define a pseudonatural equivalence from the inclusion to the constant pseudofunctor at $i$. Let $G$ be any pseudofunctor into a bicategory $\mathcal B$, and choose $\epsilon:G(i)\simeq d$. Then $\lambda_W=\epsilon G(\nu_W)$ defines a pseudonatural equivalence from $G$ restricted to $\mathcal C$ to the constant pseudofunctor at $d$. Its cell at $q$ is the actual paste
\[
 \epsilon G(\nu_V)G(q)
 \xRightarrow{\epsilon\,G_2}\epsilon G(\nu_Vq)
 \xRightarrow{\epsilon\,G(u_q)}\epsilon G(\nu_W),
\]
with associators and the chosen constant-functor units retained. The composition equation follows by the pseudofunctor pentagon for $G$ and the source equation for $u_{rq}$; the identity equation uses its unit constraint; naturality at a two-cell uses the source $u$ naturality and hom-functoriality. No local fullness, faithfulness or canonical-image premise is used. In the application, $G$ is the hom pseudofunctor of $F_m$ and $i$ is the actual identity machine. Thus the entire transported normalization, not only its object equivalences, is available.

\paragraph{Transported target normalization.} Apply $F_m$ to this already admitted native diagram, retaining its compositor and unit comparisons. Let $\epsilon_m:F_m(I_A)\simeq1_A$ be the inverse of its actual identity comparison, and put, with those frames,
\[
 \lambda_W=\epsilon_m\,F_m(\nu_W):F_m(\operatorname{idle}(W))\simeq1_A.
\]
The images of $u_q$ yield specified target three-cells
$\lambda_VF_m(q)\Rightarrow\lambda_W$ satisfying the images of the source equations. In particular $F_m(\nu_W)$ is not assumed to be a canonical unitor of $D$. It is retained as part of the normalization data. Pair and triple interface comparisons are defined using these very $\lambda$ frames and their adjunction data. When different indexing paths occur, use the already constructed individual $E_\theta$ naturality at these admitted unit-frame maps; the target joining-unit/full-witness laws provide precisely that naturality. More explicitly, the identity-compatibility modification of $E_\theta$ fixes its component at the actual identity machine, using the supplied joining/unit/assembly equations. Its native naturality three-cell at each $\nu_W$ transports that identity-machine compatibility to $\operatorname{idle}(W)$. Naturality at $u_q$ gives coherence between these transported choices. Thus this is not agreement only on objects, nor a separate arbitrary choice of interface comparisons.

\paragraph{Four-arrow equation.} The equation of perturbation pastes includes the evaluator image of the source $\Pi$. Conjugate the complete diagram by the specified $\lambda$ frames. Their adjacent factors cancel by the chosen adjunction triangles and the images of the $u_q$ composition equations. In this normalized diagram, the image of $\Pi$ is the image of the already proved source unit-frame coherence, with the actual $F_m$ constraints still present. The latter satisfy the trihomomorphism coherence equations. The remaining target identity frames satisfy the actual tricategory unit axioms of $D$. Thus the two perturbation pastes agree. This proof uses no faithful projection of $D$ and never asserts that an arbitrary image $F_m(r_0)$ or $F_m(r_2)$ is a canonical target structural map.

The identity-index/unit counterparts use the same recipe with empty fibers retained. Together with the reduction-substitution and full-witness admission checks, this supplies the explicitly listed coherent comparison $E_\theta$ with the ordered joining system. Its direction is forward semantic joining and inverse unit insertion. It is a conditional evaluation result under the stated semantic operations and their equations, not existence of those operations for every target, an equivalence of evaluator categories, or a free-algebra theorem for an unrestricted higher state monad.
\end{document}
